%%%%%%%%%%%%%%%%%%%%%%%%%%%%%%%%%%%%%%%%%%%%%%%%%%%%%%%%%%%%%%%%%%%%%%%%%%%%%%%%

\begin{abstract}
\thispagestyle{empty}
OMP4Py lleva el modelo de directivas de OpenMP al código Python. Este trabajo lo
pone a prueba adaptando cinco NAS Parallel Benchmarks de naturaleza distinta
(EP, CG, IS, MG y FT) y ejecutándolos en el supercomputador Finisterrae~III con
Python~3.15.0b1t \emph{free-threaded}.

La evaluación cubre corrección, tiempo absoluto y escalabilidad en los cuatro
modos de OMP4Py, desde el intérprete puro hasta el Cython tipado. La campaña
principal con Python~3.15.0b1t reúne 630 ejecuciones, y la comparación entre
Python~3.13.13t, 3.14.4t y 3.15.0b1t añade otras 2340, en modos interpretados y
compilados. Todas verificaron correctamente.

El factor más determinante es el tipado Cython. En clase~W, al pasar del modo
compilado sin tipos al tipado los tiempos caen en todos los benchmarks, desde el
orden de $15\times$ en FT hasta más de dos órdenes de magnitud en los casos más
favorables. La causa es directa. Sin anotaciones, cada operación del bucle
crítico arrastra el despacho de objetos Python. Con ellas, el modo~3 es la única
configuración viable para cargas grandes.

La escalabilidad entre hilos es más desigual y depende del benchmark.
EP y FT aceleran con claridad en clase~A y superan $10\times$ de speedup. Los
otros tres no escalan de forma sostenida. El diagnóstico posterior atribuye esa
diferencia a sincronización, granularidad y gestión de equipos de hilos,
combinando \texttt{perf stat} con el muestreo y trazado de Python~3.15 en
Finisterrae~III. La versión del intérprete apenas influye en los modos
compilados. En ejecución interpretada, en cambio, Python~3.13t escala bastante
peor al aumentar los hilos, salvo en EP, una penalización que el tipado elimina.
Con la afinidad de hilos corregida, Python~3.13t, 3.14t y 3.15t rinden casi igual
en EP y FT.
OMP4Py resulta útil en la práctica, pero su rendimiento no es automático y
depende del tipado Cython, de la configuración del runtime y de la estructura de
cada kernel.

\vspace*{25pt}
\begin{otherlanguage}{english}
\centerline{\bfseries \abstractname}
OMP4Py brings OpenMP's directive model to plain Python. This work puts it to the
test by adapting five NAS Parallel Benchmarks of different nature (EP, CG, IS,
MG, and FT) and running them on the Finisterrae~III supercomputer with
Python~3.15.0b1t free-threaded.

The evaluation covers correctness, absolute time, and scalability across
OMP4Py's four execution modes, from the pure interpreter to typed Cython. The
main Python~3.15.0b1t campaign gathers 630 runs, and the comparison across
Python~3.13.13t, 3.14.4t, and 3.15.0b1t adds another 2340, spanning interpreted
and compiled modes. All of them verified correctly.

By far the dominant factor is Cython typing. In class~W, moving from the untyped
compiled mode to the typed one reduces the times in every benchmark, from around
$15\times$ in FT to more than two orders of magnitude in the most favorable
cases. The reason is straightforward. Without annotations, every operation in the
critical loop drags Python object dispatch along. With them, mode~3 becomes the
only viable configuration for large workloads.

Thread scalability is more uneven and depends on the benchmark. EP and FT speed
up clearly in class~A, exceeding $10\times$ speedup. The other three do not show
sustained scalability. The later diagnosis attributes this difference to
synchronization, granularity, and thread-team management, combining
\texttt{perf stat} with Python~3.15 sampling and tracing on Finisterrae~III. The
interpreter version barely matters in the compiled modes. In interpreted
execution, by contrast, Python~3.13t scales noticeably worse as threads grow,
except on EP, a penalty that typing removes. Once thread affinity is fixed,
Python~3.13t, 3.14t, and 3.15t behave almost identically on EP and FT. OMP4Py is
useful in practice, but its performance is not automatic. It depends on Cython
typing, on the runtime configuration, and on the structure of each kernel.
\end{otherlanguage}

\vspace*{25pt}
\begin{multicols}{2}
\begin{description}
\item [Palabras chave:] \mbox{} \\
  \begin{itemize}
    \item OMP4Py
    \item OpenMP
    \item Python free-threaded
    \item NAS Parallel Benchmarks
    \item Computación de altas prestaciones
    \item Cython
    \item Escalabilidad
  \end{itemize}
\end{description}
\begin{description}
\item [Keywords:] \mbox{} \\
  \begin{itemize}
    \item OMP4Py
    \item OpenMP
    \item Free-threaded Python
    \item NAS Parallel Benchmarks
    \item High Performance Computing
    \item Cython
    \item Scalability
  \end{itemize}
\end{description}
\end{multicols}

\end{abstract}

%%%%%%%%%%%%%%%%%%%%%%%%%%%%%%%%%%%%%%%%%%%%%%%%%%%%%%%%%%%%%%%%%%%%%%%%%%%%%%%%
